\chapter{Best-Model Selection: Why Random Forest Wins for Deployment}
\label{ch:bestmodel}

The previous five chapters reported each evaluation in isolation. This
chapter aggregates the evidence and answers the practical question: if a
SOC could deploy only one of these four models tomorrow, which one
should it be?

Six criteria matter for a deployable IDS, in roughly this order of
importance: temporal stability (LODO mean and std), recall at
deployment-grade FPR budgets, calibration, inference latency,
adversarial robustness, and operational simplicity.

Table~\ref{tab:best-model} collapses every result in Part~III into a
single comparison. Each row is one criterion; each column is one model;
a star marks the best (or tied) model on that row. Random forest collects
six stars, XGBoost three, LightGBM two, and logistic regression two.

\begin{table}[H]
    \centering
    \caption{Cross-criterion comparison. $\star$ = best (or tied) on that row.}
    \label{tab:best-model}
    \footnotesize
    \begin{tabular}{lcccc}
        \toprule
        \textbf{Criterion} & \textbf{LogReg} & \textbf{Random Forest} & \textbf{XGBoost} & \textbf{LightGBM} \\
        \midrule
        Stratified multi-class macro F1 & 0.27 & 0.84 & \textbf{0.86 $\star$} & 0.30 \\
        Day-split multi-class macro F1 & 0.43 & 0.44 & 0.44 & \textbf{0.46 $\star$} \\
        Day-split binary F1 (Youden) & 0.68 & \textbf{0.86 $\star$} & 0.76 & 0.74 \\
        LODO mean ROC-AUC & 0.63 & \textbf{0.93 $\star$} & 0.92 & 0.84 \\
        LODO ROC-AUC std (lower=better) & 0.37 & \textbf{0.05 $\star$} & 0.09 & 0.21 \\
        Recall @ 0.1\% FPR (day binary) & 0.39 & \textbf{0.62 $\star$} & 0.03 & 0.46 \\
        Recall @ 1\% FPR (day binary) & 0.60 & \textbf{0.63 $\star$} & \textbf{0.64 $\star$} & 0.62 \\
        ECE (lower=better, day binary) & \textbf{0.07 $\star$} & 0.18 & 0.25 & 0.25 \\
        UNSW binary F1 & 0.89 & \textbf{0.92 $\star$} & \textbf{0.92 $\star$} & \textbf{0.92 $\star$} \\
        Latency $\mu$s/flow (bs=10k) & \textbf{0.11 $\star$} & 3.19 & 0.80 & 2.80 \\
        Adversarial evasion @ $0.25\sigma$ & n/a & 21.4\% & n/a & n/a \\
        \midrule
        \textbf{Total $\star$} & 2 & \textbf{6 $\star$} & 3 & 2 \\
        \bottomrule
    \end{tabular}
\end{table}

\interpretation{Random forest dominates the rows that matter most for
deployment: temporal stability, recall at strict FPR, and overall
ranking across LODO folds. It loses on stratified multi-class (XGBoost
wins) and on calibration (LogReg wins). For deployment, the wins it has
are the ones that matter.}

Random forest is the best deployment choice on three independent
grounds. \textbf{Temporal stability}: ROC-AUC $= 0.926 \pm 0.053$ across
LODO folds is the lowest standard deviation of any of the four models
(XGBoost $\pm 0.086$, LightGBM $\pm 0.212$, LogReg $\pm 0.374$).
Stability matters more than peak performance because production traffic
shifts continuously. \textbf{Operational sweet spot at strict FPR}: at
0.1\% FPR the random forest catches 62\% of attacks, while XGBoost
catches 3\%, LightGBM catches 46\%, and logistic regression catches
39\%. \textbf{Reasonable speed and simplicity}: random forest is the
slowest of the four at 3.2~$\mu$s/flow, but still 200 times faster than
the throughput a saturated 1~Gbps mirror tap produces. Its single-step
training and native parallel evaluation make it operationally simpler
than gradient-boosting alternatives.

Three caveats apply. Multi-class on the day split is hard for random
forest just as for the other tree ensembles --- all three reach macro F1
of approximately 0.44 because of the unseen-class problem. Calibration
on the day split is poor (ECE 0.18 versus LogReg's 0.07), and any
deployment that exposes probability scores to downstream tooling should
apply Platt scaling or isotonic regression as post-processing. Finally,
adversarial robustness is shared with the other tree models: the 21.4\%
evasion rate at $\varepsilon = 0.25\sigma$ (or 24.4\% if the attacker is
allowed to run unbounded) is not made better by choosing random forest
over XGBoost or LightGBM, and network-layer mitigations (rate limiting,
traffic shaping, or protocol-level constraints on inter-packet timing)
are needed in addition to the model itself; their effectiveness would
need packet-level validation outside this thesis's scope.

The deployment recommendation: random forest as the binary detector
with Youden-J threshold tuning per fold, retrained weekly on a rolling
four-week window. Keep XGBoost as a transferability sanity check rather
than as a defence: adversarial flows that evade random forest transfer
to XGBoost at 86\%, so the two models do not catch independent failure
modes, but a divergence between them on a particular flow is a useful
out-of-distribution signal. Reserve LightGBM for UNSW-style multi-class
tasks where it leads. Use logistic regression as the calibration
baseline.
